\section{Intents: MCP Tool Generation}
\label{sec:intents}

Where an action exposes a REST route or CLI command, an
\texttt{intent} exposes the same kind of typed signature to an AI
agent: a named, model-callable tool with a declared input and output
shape, compiled into an MCP (Model Context Protocol) server's
\texttt{tools/list} the way an action compiles into an HTTP route.
Intents live in their own top-level \texttt{intents} list, sibling to
\texttt{actions}.

\begin{lstlisting}
name: gomcp
intents:
  - name: add
    title: Add two numbers
    description: Adds two numbers and returns their sum.
    annotations: {readOnlyHint: true, idempotentHint: true}
    in:
      fields:
        - {name: a, type: float64}
        - {name: b, type: float64}
    out:
      fields: [{name: sum, type: float64}]

  - name: deleteInvoiceTool
    from: deleteInvoice
    title: Delete an invoice
    annotations: {destructiveHint: true}

actions:
  - name: searchInvoices
    method: get
    url: /invoices/search
    intent: true
    in:
      fields: [{name: query, type: string}]
    out:
      fields: [{name: total, type: int}]
\end{lstlisting}

This one file shows all three ways to declare an intent. \texttt{add}
is standalone: no backing action exists, since a pure math tool has no
reason to be an HTTP route. \texttt{deleteInvoiceTool} is derived from
an existing action via \texttt{from} --- \texttt{description}/\texttt{in}/
\texttt{out} are copied from \texttt{deleteInvoice} wherever the
intent doesn't declare its own, so exposing an action as a tool needs
no field duplication. \texttt{searchInvoices} needs no \texttt{intents}
entry at all: \texttt{intent: true} on the action synthesizes one
automatically, under the action's own name.

\begin{sidebar}{Intent properties}
\small
\begin{tabular}{@{}p{0.24\linewidth}p{0.7\linewidth}@{}}
\textbf{name} & tool name shown to the model; defaults to \texttt{from} \\
\textbf{from} & name of an existing action to derive description/in/out from \\
\textbf{title} & human-readable title (MCP's \texttt{title} annotation) \\
\textbf{description} & sent verbatim to the model --- write it for the model, not for a developer \\
\textbf{in} / \textbf{out} & request/response body, same shape as an action's \\
\textbf{annotations} & \texttt{readOnlyHint}, \texttt{destructiveHint}, \texttt{idempotentHint}, \texttt{openWorldHint} \\
\textbf{permission} & key from the module's permission tree gating visibility \\
\end{tabular}
\end{sidebar}

Preprocessing resolves \texttt{intent: true} and \texttt{from}
references before any compiler runs, and validates that a
\texttt{dto} reference on an intent's \texttt{in}/\texttt{out} names a
dto the module actually declares --- a typo surfaces as a clear error
at compile time instead of an ``undefined type'' failure deep inside
generated Go.

\subsection{What the Go compiler generates}
\texttt{emi go} renders every intent into \texttt{Intents.go}: an
\texttt{\{Name\}IntentArgs}/\texttt{\{Name\}IntentResult} struct pair
for any intent with inline fields (a \texttt{dto} or \texttt{primitive}
reference reuses the existing type instead, exactly as an action
body would), through the very same struct generator dtos and actions
use --- so nested objects, arrays, enums, \texttt{complex} fields and
CLI flag helpers all work identically. Alongside each struct pair sits
an accessor carrying the tool's metadata:

\begin{lstlisting}
func AddIntentMeta() struct {
    Name, Title, Description, Permission string
    ReadOnlyHint, DestructiveHint,
    IdempotentHint, OpenWorldHint bool
}
\end{lstlisting}

This is deliberately as far as generation goes: wiring
\texttt{\{Name\}IntentMeta()} and the typed \texttt{Args}/\texttt{Result}
structs into a real \texttt{*mcp.Server} is left to a hand-written file,
the same division of labor a reactive action gets between its
generated file and a developer-supplied \texttt{Impl} --- \texttt{emi}
has no opinion on which MCP SDK a project standardizes on, or on the
actual tool implementation. \texttt{examples/go-mcp} shows the whole
loop end to end: a two-tool math module compiled with \texttt{emi go},
served for real over MCP's stdio JSON-RPC transport by a few dozen
lines of hand-written Go.

\begin{pullquote}
Because \texttt{intents} reuses the same field/struct/complex-type
machinery as actions and dtos, an intent is never a second definition
of a shape that already exists elsewhere in the module --- \texttt{from}
and \texttt{dto} both mean ``point at the type that's already there.''
\end{pullquote}
